% sesmesh — Solvent Excluded Surface meshes for molecules
% Copyright (c) 2026 Ugo Varetto
% SPDX-License-Identifier: MIT OR Apache-2.0

\section{Implementation}
\label{sec:implementation}

\sesmesh{} is written in Rust (edition 2024, minimum supported compiler
1.89) as a Cargo workspace of a library crate, \code{sesmesh}, a
command-line crate, \code{sesmesh-cli} (binary \code{sesmesh}), and a
development-only crate of test oracles, analytic references and cluster
generators. All crates forbid \code{unsafe} code, every public item is
documented, the library never prints (warnings are returned as values), and
the public API exposes no third-party types, so dependency upgrades cannot
break callers. Since 1.0 the project follows semantic versioning, and the
1.x series keeps the default output byte-identical; new behaviour is
opt-in.

\subsection{Library and command-line interface}

The core entry point takes spheres and parameters and returns a mesh with
per-vertex normals, nearest atoms, patch types and components, plus
statistics; structure readers, atom selection and radius assignment are
separate composable functions (\cref{lst:api}). A reusable
\code{SurfaceEngine} keeps per-thread buffers between calls for
trajectories, and a \code{Control} object carries a cancellation flag and a
progress callback; cancellation is checked per sphere, between stages and
per brick, and takes effect within one brick's work (under 100\,ms in the
tests). The command-line tool exposes every library option, writes one or
many inputs (batch mode with per-file error isolation), reports timings and
statistics, generates shell completions and a manual page, and uses
distinct exit codes (0 success, 1 runtime error, 2 usage error, 3 input
parse error, 4 resource limit, 130 interrupted).

\begin{figure}[t]
\begin{lstlisting}[language=Rust]
use sesmesh::export::obj::{write_obj_path, ObjOptions};
use sesmesh::{compute_surface, io, radii, Diagnostics, Selection, Sphere, SurfaceParams};

fn main() -> sesmesh::Result<()> {
    let structure = io::read_path("1crn.cif", &io::ReadOptions::default())?;
    let mut diags = Diagnostics::new();
    let atoms = structure.select(&Selection::default(), &mut diags)?;
    let radii = radii::assign(&atoms, &radii::RadiusOptions::default(), &mut diags)?;
    let spheres = Sphere::from_atoms(&atoms, &radii)?;
    let out = compute_surface(&spheres, &SurfaceParams::ses(1.4).with_spacing(0.25))?;
    println!("{} triangles, area {:.1} A^2", out.stats.triangles, out.stats.area);
    write_obj_path("1crn.obj", &out.mesh, &ObjOptions::default())
}
\end{lstlisting}
\begin{lstlisting}
sesmesh 1crn.cif -q high --cavities remove --vertex-colors element --stats stats.json
sesmesh 1k4r.cif.gz --assembly 1 -q draft -o capsid.ply      # 545k-atom capsid
sesmesh 3j3q.cif.gz --max-memory 4G --auto-coarsen          # coarsen until it fits
\end{lstlisting}
  \caption{Library use (from the README; the snippets on the project
    website are compiled against 1.1.0) and three command-line
    invocations.}
  \label{lst:api}
\end{figure}

\paragraph{Inputs.}
Purpose-built readers parse PDB (fixed columns, hybrid-36 serial
numbers~\cite{pdbformat33,hybrid36}, chunked parallel parsing of large
files), PDBx/mmCIF~\cite{mmcif} (a full CIF~1.1 lexer and the
\code{\_atom\_site} loop in any column order; about 350\,MiB/s on one core;
the gzipped 2.44-million-atom 3J3Q entry reads in about 0.8\,s), XYZ
and extended XYZ~\cite{extxyz}, and PQR~\cite{pqr}, each optionally gzip
compressed~\cite{flate2}. Biological assemblies are generated from mmCIF
\code{pdbx\_struct\_\allowbreak assembly\_gen} and
\code{pdbx\_struct\_\allowbreak oper\_list} and PDB
\code{REMARK 350}; the 1K4R entry contains 9\,084 deposited atoms of a
capsid of 545\,040. Assemblies of more than 100\,000 copies or 500 million
atoms and decompressed inputs above 8\,GiB are refused, which bounds memory
on hostile input. The parsers are tested by reading the same entries in
PDB and mmCIF format, which must give identical atom sets, and against the
deposited atom counts published by the RCSB. Selection defaults to the first model, the
highest-occupancy alternate location per residue, and no hydrogens or
waters. Radius sets are Bondi~\cite{bondi1964} completed with Mantina et
al.~\cite{mantina2009} (default), Alvarez~\cite{alvarez2013}, the
united-atom ProtOr radii~\cite{tsai1999}, a uniform radius, per-file radii
(PQR, extended XYZ) and custom TOML rules.

\paragraph{Outputs.}
Wavefront OBJ~\cite{objformat} with a header recording the version, input,
every parameter and summary statistics, written by a parallel fixed-point
formatter (about 23\,ms for the 1.28-million-triangle 4HHB mesh), with
optional per-component groups and per-vertex colours by element, chain,
patch type or component; binary PLY with per-vertex atom, patch and
component attributes; and the MSMS \code{.vert}/\code{.face}
format~\cite{msmsmanual} for programs that consume MSMS meshes. All three
can be gzip compressed or written to standard output. Side files give
per-vertex attributes and per-atom SES areas (each triangle's area split
equally among the nearest atoms of its vertices) as CSV, and statistics as
JSON.

\subsection{Parallelism and determinism}
\label{sec:determinism}

Parallelism uses rayon's work stealing~\cite{rayon}. The library never
creates a global thread pool: it runs in the caller's pool, and the CLI
builds one from \code{--threads}. Parsing, radius assignment, cell lists,
the arrangement, the touched set, classification, the fused per-brick
pipeline (gather, kernel, Marching Cubes, refinement), prefix sums,
reference resolution, components, validation, smoothing and output
formatting are parallel. Without the \code{parallel} feature the same code
runs sequentially.

Output is \emph{byte-identical} for identical input, options, binary and
SIMD level, independent of the number of threads and of scheduling. Four
rules achieve this: per-item results are collected into indexed vectors and
concatenated in order; every floating-point total is a sum of per-brick
(or per-chunk) partials added sequentially in a fixed order; union--find
links the larger root to the smaller (\cref{sec:cavities}); and parallel
sorts break ties by index. Memory-saving chunking
($\mathrm{clamp}(n/256, 1, 4096)$ items per chunk) is likewise invisible in
the result. A test suite compares OBJ files byte for byte for 1, 2, 4, 8
and 16 threads and for repeated runs on crambin and trypsin (2PTN) in
continuous integration, and on 4HHB and 1AON nightly. Determinism also
makes memory-limited runs reproducible: the memory estimate deliberately
ignores the thread count, so the automatically coarsened spacing, and with
it the mesh, cannot depend on it (\cref{sec:memory}).

\subsection{SIMD with runtime dispatch}

Stable Rust has no portable SIMD in its standard library~\cite{portablesimd};
\sesmesh{} uses the safe \code{fearless\_simd} crate~\cite{fearlesssimd},
which detects the CPU at run time and dispatches generic kernels at SSE4.2,
AVX2, AVX-512 or NEON width, with a portable fallback. The level is
detected once per engine and can be forced for diagnostics; an explicit
level is capped at the detected one, and a level of another architecture
is a usage error. Sampled field values at every available level agree with
the fallback within $2\cdot10^{-5}$\Ang{} (a test run on x86-64 and arm64
runners in continuous integration). On Apple Silicon the compiler
auto-vectorizes the fallback too, and NEON was only 3--6\,\% faster than it
(crambin and haemoglobin at 0.5 and 0.25\Ang); the AVX2 versus SSE2 ratio
on x86-64 has not been measured yet (\cref{sec:discussion}). Marching Cubes
case classification stays scalar (about 2\,\% of a single-thread profile).

\subsection{Memory model and memory guard}
\label{sec:memory}

Peak memory has a per-atom part (the arrangement and the structure) and a
per-triangle part (brick meshes, assembly, the final mesh). Compacting the
arrangement---neighbour rows only for exposed spheres, no stored arcs,
circles without their second basis vector, exact capacities---reduced the
field data on the 1K4R capsid from about 1\,150 to 560 bytes per atom, plus
about 150 bytes per atom for the input structure and spheres; building it
in waves of flat per-chunk arrays and freeing the field before assembly
lowered peaks further (\cref{sec:large}).

A memory guard (\code{--max-memory}) predicts the peak before meshing. It
counts the Marching Cubes triangles of all surface bricks, without creating
vertices, when there are at most 512, and otherwise of every $k$-th brick
in key order, extrapolating $T = n(\bar t + 2 s_t/\sqrt m)$ from the mean
and standard deviation over the $m$ sampled bricks. For grids finer than
the pilot spacing (brick side 5\Ang), it classifies the pilot grid and
predicts $T \propto h^{-2}$ (times 1.1) and the number of touched bricks
$X \propto h^{-3}$, so it never classifies a grid that would not fit. The
estimate is
\begin{equation}
  M = 16\,\text{MiB} + 1200\,N + (150 + 50\,[\text{decimation}]
    + 50\,[\text{smoothing}])\,T + 100\,X \ \text{bytes},
\end{equation}
calibrated on measured peaks; on the test runs it was 1.15--1.7 times the
measured peak resident set size. Above the limit the run is refused (exit
code 4), or, with \code{--auto-coarsen}, the smallest spacing whose
prediction fits is found by bisection and checked on its own sample, and if
needed $h$ grows by $1.05\sqrt{(M - M_0)/(L - M_0)}$ ($L$ the limit, $M_0$
the part independent of $h$), at most four times. For example, the 1K4R
capsid requested at 0.25\Ang{} under a 1\,GiB limit is meshed at 2.08\Ang{}
with an estimate of 1\,001\,MiB and a measured peak of 600\,MiB, and 1AON
requested at 0.1\Ang{} under 512\,MiB at 0.634\Ang{} (estimate 497\,MiB,
peak 352\,MiB).

\subsection{Quality assurance}

Every optimized path has a simple reference: a brute-force double-precision
oracle of $F$ over all features, against which the exact point evaluator,
the SIMD kernel at every level and the lattice samples of random clusters
are checked (property-based tests with random radii and probes); analytic
surfaces (\cref{sec:analytic}); the exhaustive Marching Cubes table tests
(\cref{sec:mc}); mesh validation of every result; and the determinism suite.
Seven fuzz targets cover the PDB, mmCIF, CIF lexer, XYZ, PQR, format
detection and assembly code; before 1.0 each ran for 10 minutes (about 18
million inputs in total) without a crash, timeout or out-of-memory report,
and a longer campaign is pending (\cref{sec:discussion}). Continuous
integration runs formatting, a strict lint configuration, tests on macOS
arm64, Linux x86-64, Linux arm64 and Windows x86-64, documentation, licence
and advisory checks, the minimum compiler version, and five runnable
examples; a nightly job runs the long tests on downloaded structures.
